%\TODO{We should mention here and earlier that we retrieve from a subset of the websites (approved by wikimedia).}
% \yucheng{yep. I added a footnote in section 3.4 and also revised the ethics statement with this information.}

We build and evaluate our work to strictly adhere to ethical standards. The construction of the WildSeek dataset involves collecting data with users’ explicit approval, and we carefully remove all personally identifiable information. In contrast to creative generation tasks, our tasks generate content that may impact how people perceive information and shape their opinions. We design our system to ground generated content on openly accessible external sources available on the general internet, with proper citations. Our experiments and evaluations ensure the accurate delivery of information and significantly reduce hallucinations. We avoid publishing or posting any generated content without careful examination of information accuracy. We believe there are no data privacy issues as we ground our generated content from information accessible to the general public.

The primary risk of our work is the common bias issues originating from biases present on the general internet. Following \citet{shao2024assisting}, we mitigate this problem by applying rule-based filtering according to the Wikipedia guideline\footnote{\url{https://en.wikipedia.org/wiki/Wikipedia:Reliable_sources}} and incorporating multiple sources. Additionally, in our web application, we have implemented input topic moderation to reject topics that are sensitive, illegal, or potentially violate personal privacy. However, further information processing modules that serve as filters for internet sources and more robust modules to verify the accuracy of information can be implemented. Additionally, our current work only considers generating and retrieving information from English sources. Extending our system to be compatible with multilingual sources and generation will be beneficial.